
\section{Advanced Readiness Activities}

The following activities are intended for strong senior undergraduate,
master's, or beginning PhD students. They reinforce the mathematical and
computational foundations needed for the later control co-design chapters.
The problems combine multivariable calculus, numerical integration, linear
algebra, ordinary differential equations, state-space modeling, discretization,
and basic optimization.

\subsection*{Problem 1: Composite Engineering Functions, Gradients, and Hessians}

For a mass--spring--damper system, define the natural frequency and damping
ratio as
\begin{equation}
\omega_n(m,k)
=
\sqrt{\frac{k}{m}},
\end{equation}
and
\begin{equation}
\zeta(m,k,c)
=
\frac{c}{2\sqrt{km}}.
\end{equation}

Consider the dimensionless design objective
\begin{equation}
J(m,k,c)
=
\left(
\frac{\omega_n-\omega_d}{\omega_d}
\right)^2
+
\beta
\left(
\frac{\zeta-\zeta_d}{\zeta_d}
\right)^2
+
\gamma
\left(
\frac{m}{m_{\mathrm{ref}}}
\right)^2,
\end{equation}
where
\begin{equation}
\omega_d=4~\mathrm{rad/s},
\qquad
\zeta_d=0.7,
\qquad
\beta=2,
\end{equation}
\begin{equation}
\gamma=0.05,
\qquad
m_{\mathrm{ref}}=2~\mathrm{kg}.
\end{equation}

\begin{enumerate}
    \item Derive
    \begin{equation}
    \frac{\partial\omega_n}{\partial m},
    \qquad
    \frac{\partial\omega_n}{\partial k}.
    \end{equation}

    \item Derive
    \begin{equation}
    \frac{\partial\zeta}{\partial m},
    \qquad
    \frac{\partial\zeta}{\partial k},
    \qquad
    \frac{\partial\zeta}{\partial c}.
    \end{equation}

    \item Use the chain rule to derive the complete gradient
    \begin{equation}
    \nabla J(m,k,c).
    \end{equation}

    \item Evaluate $J$ and $\nabla J$ at
    \begin{equation}
    m=2~\mathrm{kg},
    \qquad
    k=24~\mathrm{N/m},
    \qquad
    c=5~\mathrm{N\,s/m}.
    \end{equation}

    \item Let
    \begin{equation}
    \mathbf{d}
    =
    \frac{1}{\sqrt{6}}
    \begin{bmatrix}
    1\\
    -2\\
    1
    \end{bmatrix}.
    \end{equation}
    Compute the directional derivative
    \begin{equation}
    D_{\mathbf{d}}J
    =
    \nabla J^T\mathbf{d}.
    \end{equation}

    \item Derive the $3\times3$ Hessian matrix
    \begin{equation}
    H_J(m,k,c).
    \end{equation}

    \item Evaluate the eigenvalues of the Hessian at the stated design and
    determine whether $J$ is locally convex there.

    \item Verify the analytical gradient using central finite differences
    with step sizes
    \begin{equation}
    h=10^{-2},\ 10^{-4},\ 10^{-6},\ 10^{-8}.
    \end{equation}
\end{enumerate}


\subsection*{Problem 2: Numerical Integration and Observed Convergence Order}

Consider
\begin{equation}
q(t)
=
e^{-t}\sin(3t),
\qquad
0\leq t\leq2,
\end{equation}
and define
\begin{equation}
I
=
\int_0^2 q(t)^2\,dt.
\end{equation}

\begin{enumerate}
    \item Evaluate $I$ analytically.

    \item Approximate $I$ using:
    \begin{enumerate}
        \item the left-rectangle rule,
        \item the midpoint rule,
        \item the composite trapezoidal rule,
        \item composite Simpson's rule.
    \end{enumerate}

    \item Use
    \begin{equation}
    N=10,\ 20,\ 40,\ 80,\ 160
    \end{equation}
    equal subintervals.

    \item For each method, compute the absolute error
    \begin{equation}
    e_N
    =
    |I_N-I|.
    \end{equation}

    \item Estimate the observed convergence order using
    \begin{equation}
    p_N
    =
    \frac{
    \log(e_N/e_{2N})
    }{
    \log 2
    }.
    \end{equation}

    \item Compare the observed orders with the theoretical orders of the four
    methods.

    \item Derive the composite trapezoidal approximation in matrix form:
    \begin{equation}
    I_N
    =
    \mathbf{w}^T\mathbf{q}^{\,2},
    \end{equation}
    where $\mathbf{w}$ is the quadrature-weight vector.

    \item Explain why a highly accurate state trajectory does not
    automatically guarantee a highly accurate integral performance measure
    if the quadrature rule is inconsistent with the trajectory
    approximation.
\end{enumerate}


\subsection*{Problem 3: Linear Algebra, Modal Coordinates, and Positive Definiteness}

Consider the undamped two-degree-of-freedom mechanical system
\begin{equation}
M\ddot{\mathbf{q}}
+
K\mathbf{q}
=
\mathbf{0},
\end{equation}
where
\begin{equation}
M=
\begin{bmatrix}
2&0\\
0&1
\end{bmatrix},
\qquad
K=
\begin{bmatrix}
6&-2\\
-2&4
\end{bmatrix}.
\end{equation}

\begin{enumerate}
    \item Verify that $M$ and $K$ are symmetric positive definite.

    \item Solve the generalized eigenvalue problem
    \begin{equation}
    K\boldsymbol{\phi}
    =
    \omega^2M\boldsymbol{\phi}.
    \end{equation}

    \item Compute the two natural frequencies.

    \item Normalize the mode shapes so that
    \begin{equation}
    \Phi^TM\Phi=I.
    \end{equation}

    \item Verify that
    \begin{equation}
    \Phi^TK\Phi
    =
    \operatorname{diag}
    \left(
    \omega_1^2,
    \omega_2^2
    \right).
    \end{equation}

    \item Introduce modal coordinates
    \begin{equation}
    \mathbf{q}=\Phi\boldsymbol{\eta}
    \end{equation}
    and derive the decoupled equations of motion.

    \item For
    \begin{equation}
    \mathbf{q}(0)=
    \begin{bmatrix}
    1\\
    0
    \end{bmatrix},
    \qquad
    \dot{\mathbf{q}}(0)=
    \begin{bmatrix}
    0\\
    0
    \end{bmatrix},
    \end{equation}
    compute the initial modal coordinates.

    \item Derive the exact displacement response
    \begin{equation}
    \mathbf{q}(t).
    \end{equation}

    \item Compute the total mechanical energy
    \begin{equation}
    E(t)
    =
    \frac{1}{2}
    \dot{\mathbf{q}}^TM\dot{\mathbf{q}}
    +
    \frac{1}{2}
    \mathbf{q}^TK\mathbf{q}
    \end{equation}
    and prove that it is constant.

    \item Explain why positive definiteness of $M$ and $K$ matters
    physically and mathematically.
\end{enumerate}


\subsection*{Problem 4: Exact State-Space Response of a Forced Dynamic System}

Consider
\begin{equation}
m\ddot{x}
+
c\dot{x}
+
kx
=
u_0,
\end{equation}
where $u_0$ is a constant input. Use
\begin{equation}
m=1.5~\mathrm{kg},
\qquad
c=1.8~\mathrm{N\,s/m},
\qquad
k=12~\mathrm{N/m},
\end{equation}
\begin{equation}
u_0=3~\mathrm{N},
\qquad
x(0)=0.2~\mathrm{m},
\qquad
\dot{x}(0)=-0.1~\mathrm{m/s}.
\end{equation}

\begin{enumerate}
    \item Define
    \begin{equation}
    \mathbf{x}
    =
    \begin{bmatrix}
    x\\
    \dot{x}
    \end{bmatrix}
    \end{equation}
    and derive
    \begin{equation}
    \dot{\mathbf{x}}
    =
    A\mathbf{x}
    +
    B u_0.
    \end{equation}

    \item Compute the eigenvalues of $A$ and classify the response as
    underdamped, critically damped, or overdamped.

    \item Compute the equilibrium state
    \begin{equation}
    \mathbf{x}_{\mathrm{eq}}
    \end{equation}
    associated with the constant input.

    \item Define
    \begin{equation}
    \widetilde{\mathbf{x}}
    =
    \mathbf{x}
    -
    \mathbf{x}_{\mathrm{eq}}
    \end{equation}
    and derive the homogeneous shifted system.

    \item Derive
    \begin{equation}
    \mathbf{x}(t)
    =
    e^{At}
    \left(
    \mathbf{x}(0)-\mathbf{x}_{\mathrm{eq}}
    \right)
    +
    \mathbf{x}_{\mathrm{eq}}.
    \end{equation}

    \item Obtain explicit scalar expressions for $x(t)$ and $\dot{x}(t)$.

    \item Compute the first time at which the displacement reaches its
    steady-state value.

    \item Compute the maximum displacement over
    \begin{equation}
    0\leq t\leq5.
    \end{equation}

    \item Verify the analytical solution using a numerical ODE integrator.

    \item Repeat the numerical solution using forward Euler with
    \begin{equation}
    h=0.2,\ 0.1,\ 0.05,\ 0.025.
    \end{equation}
    Estimate the observed global convergence order.
\end{enumerate}


\subsection*{Problem 5: Discretized Minimum-Energy Control Solved by Linear Algebra}

Consider the double-integrator system
\begin{equation}
\dot{x}_1=x_2,
\qquad
\dot{x}_2=u,
\end{equation}
with
\begin{equation}
x_1(0)=0,
\qquad
x_2(0)=0.
\end{equation}

Use a time horizon
\begin{equation}
T=1
\end{equation}
and divide it into four equal intervals of length
\begin{equation}
h=\frac{1}{4}.
\end{equation}

Assume the control is piecewise constant:
\begin{equation}
u(t)=u_i,
\qquad
t\in[ih,(i+1)h),
\qquad
i=0,1,2,3.
\end{equation}

Use the exact zero-order-hold discrete dynamics
\begin{equation}
\begin{bmatrix}
x_{1,i+1}\\
x_{2,i+1}
\end{bmatrix}
=
\begin{bmatrix}
1&h\\
0&1
\end{bmatrix}
\begin{bmatrix}
x_{1,i}\\
x_{2,i}
\end{bmatrix}
+
\begin{bmatrix}
h^2/2\\
h
\end{bmatrix}
u_i.
\end{equation}

The terminal requirements are
\begin{equation}
x_1(1)=1,
\qquad
x_2(1)=0.
\end{equation}

Minimize the discrete control effort
\begin{equation}
J
=
h\sum_{i=0}^{3}u_i^2.
\end{equation}

\begin{enumerate}
    \item Propagate the dynamics symbolically and express the terminal state
    as
    \begin{equation}
    \mathbf{x}_4
    =
    G\mathbf{U},
    \end{equation}
    where
    \begin{equation}
    \mathbf{U}
    =
    \begin{bmatrix}
    u_0&u_1&u_2&u_3
    \end{bmatrix}^{T}.
    \end{equation}

    \item Write the terminal requirements as
    \begin{equation}
    G\mathbf{U}
    =
    \mathbf{b}.
    \end{equation}

    \item Show that the problem is a minimum-norm equality-constrained
    quadratic problem.

    \item Using Lagrange multipliers, prove that
    \begin{equation}
    \mathbf{U}^*
    =
    G^T
    \left(
    GG^T
    \right)^{-1}
    \mathbf{b}.
    \end{equation}

    \item Compute all four optimal control values.

    \item Compute the complete state trajectory at the five time nodes.

    \item Verify the terminal conditions exactly.

    \item Repeat the derivation for an arbitrary number $N$ of equal
    intervals.

    \item Explain how the dimensions of $G$, $\mathbf{U}$, and
    $\mathbf{b}$ change with $N$.

    \item Compare the discrete solution with the continuous minimum-energy
    control obtained by calculus of variations or an optimal-control
    solver.
\end{enumerate}


\subsection*{Problem 6: State and Parameter Sensitivities for a Readiness Mini-Project}

Consider the mass--spring--damper system
\begin{equation}
m\ddot{q}
+
c\dot{q}
+
kq
=
u(t),
\end{equation}
with
\begin{equation}
m=1,
\qquad
q(0)=1,
\qquad
\dot{q}(0)=0,
\end{equation}
and
\begin{equation}
u(t)=e^{-t}.
\end{equation}

The plant parameters are
\begin{equation}
\mathbf{p}
=
\begin{bmatrix}
k\\
c
\end{bmatrix}.
\end{equation}

Define the state vector
\begin{equation}
\mathbf{x}
=
\begin{bmatrix}
q\\
v
\end{bmatrix},
\end{equation}
and the performance measure
\begin{equation}
J(k,c)
=
\frac{1}{2}
\int_0^5
\left(
q(t)^2
+
0.1v(t)^2
\right)dt
+
0.01k^2
+
0.02c^2.
\end{equation}

\begin{enumerate}
    \item Derive the state-space model
    \begin{equation}
    \dot{\mathbf{x}}
    =
    \mathbf{f}
    \left(
    \mathbf{x},
    \mathbf{p},
    t
    \right).
    \end{equation}

    \item Define the parameter-sensitivity vectors
    \begin{equation}
    \mathbf{s}_k
    =
    \frac{\partial\mathbf{x}}{\partial k},
    \qquad
    \mathbf{s}_c
    =
    \frac{\partial\mathbf{x}}{\partial c}.
    \end{equation}

    \item Derive the sensitivity equations
    \begin{equation}
    \dot{\mathbf{s}}_j
    =
    \frac{\partial\mathbf{f}}{\partial\mathbf{x}}
    \mathbf{s}_j
    +
    \frac{\partial\mathbf{f}}{\partial p_j},
    \qquad
    j\in\{k,c\}.
    \end{equation}

    \item Derive the initial conditions for both sensitivity vectors.

    \item Derive
    \begin{equation}
    \frac{\partial J}{\partial k}
    \qquad\text{and}\qquad
    \frac{\partial J}{\partial c}
    \end{equation}
    in terms of the state and sensitivity trajectories.

    \item Simulate the state and sensitivity equations simultaneously at
    \begin{equation}
    k=4,
    \qquad
    c=0.8.
    \end{equation}

    \item Verify both derivatives using central finite differences.

    \item Investigate the finite-difference error for
    \begin{equation}
    h=10^{-2},\ 10^{-4},\ 10^{-6},\ 10^{-8}.
    \end{equation}

    \item Use the computed gradient to perform ten steepest-descent
    iterations with a backtracking line search, subject to
    \begin{equation}
    0.5\leq k\leq10,
    \qquad
    0.05\leq c\leq4.
    \end{equation}

    \item Report the final design, objective value, gradient norm, and state
    response.

    \item Explain how this mini-project combines the main prerequisite topics
    of Chapter 0 and prepares students for later CCD formulations.
\end{enumerate}
